% BHU PYQ -- Auto-generated & error-corrected
\documentclass[11pt,a4paper]{article}

\usepackage[a4paper,top=2.5cm,bottom=2.5cm,left=2.8cm,right=2.8cm,headheight=14pt]{geometry}
\usepackage{amsmath,amssymb}
\usepackage[T1]{fontenc}
\usepackage[utf8]{inputenc}
\usepackage{lmodern}
\usepackage{microtype}
\usepackage{fancyhdr}
\usepackage{enumitem}
\usepackage{hyperref}
\usepackage{lastpage}
\usepackage{parskip}

\hypersetup{colorlinks=true,urlcolor=black,linkcolor=black,
  pdftitle={STB-101: Statistical Methods and Probability -- BHU PYQ},pdfauthor={Banaras Hindu University}}

\pagestyle{fancy}\fancyhf{}
\renewcommand{\headrulewidth}{0.4pt}\renewcommand{\footrulewidth}{0.4pt}
\fancyhead[L]{\small\textit{Banaras Hindu University}}
\fancyhead[R]{\small\textit{STB-101: Statistical Methods and Probability}}
\fancyfoot[C]{\small Page \thepage\ of \pageref{LastPage}}

\newlist{parts}{enumerate}{2}
\setlist[parts,1]{label=(\alph*),leftmargin=*,itemsep=5pt,topsep=3pt}
\setlist[parts,2]{label=(\roman*),leftmargin=*,itemsep=3pt,topsep=2pt}
\newcommand{\pts}[1]{\hfill\textit{[\#1]}}

\begin{document}

\begin{center}
  {\large\bfseries BANARAS HINDU UNIVERSITY}\\[2pt]
  {\normalsize B.Sc. (Hons.) Semester I Examination, 2022-23}\\[6pt]
  \rule{0.8\linewidth}{0.5pt}\\[6pt]
  {\large\bfseries Paper No. STB-101: Statistical Methods and Probability}\\[4pt]
  {\normalsize Time: 3 Hours\hfill Maximum Marks: 70}
\end{center}
\rule{\linewidth}{0.4pt}
\smallskip
\noindent\textit{Instructions: Answer \textbf{five} questions including Question~1 which is compulsory. Symbols have their usual meanings.}
\bigskip

\medskip\noindent\textbf{Question 1.}
\begin{parts}
  \item What do you mean by measures of central tendency? Also, describe the procedures of construction of a frequency curve, frequency polygon, and histogram.
  \item Define geometric mean. Also, write down its merits and demerits. Further, let $G_1$ and $G_2$ be the geometric means of two individual series and $n_1$ and $n_2$ be their sizes respectively, then derive the formula for the geometric mean of the combined series.
\end{parts}
\medskip\rule{\linewidth}{0.2pt}

\medskip\noindent\textbf{Question 2.}
\begin{parts}
  \item Define moments about origin and moments about arithmetic mean. Further, derive the relation between moments about mean in terms of moments about origin.
  \item Define mean square deviation and variance. Also, calculate the mean deviation from the mean and standard deviation of the series $a, a+d, a+2d, \dots, a+2nd$ and verify that the latter is greater than the former.
\end{parts}
\medskip\rule{\linewidth}{0.2pt}

\medskip\noindent\textbf{Question 3.}
\begin{parts}
  \item Give the axiomatic definition of probability. Next, if $A$ and $B$ be any two events of a random experiment, then find the probability of happening of at least $A$ or $B$.
  \item If $a$ and $b$ be any two non-negative quantities such that $a+b=2n$, then find the probability that their product is not less than $\frac{3}{4}$ times their greatest product.
\end{parts}
\medskip\rule{\linewidth}{0.2pt}

\medskip\noindent\textbf{Question 4.}
\begin{parts}
  \item State and prove a necessary and sufficient condition for any two events $A$ and $B$ of a random experiment such that $P(A)>0$ and $P(B)>0$ to be independent events.
  \item Define conditional probability. Further, if a certain drug manufactured by a company is tested chemically for its toxic nature. Let the event ``the drug is toxic'' be denoted by $E$ and the event ``the chemical test reveals that the drug is toxic'' be denoted by $F$. Let $P(E) = \lambda$, $P(F|E) = 1 - \lambda$, and $P(F^c|E^c) = 1 - \lambda$. Then prove that the probability that the drug is not toxic given that the chemical test reveals that it is toxic, is independent of $\lambda$.
\end{parts}
\medskip\rule{\linewidth}{0.2pt}

\medskip\noindent\textbf{Question 5.}
\begin{parts}
  \item If $E_1, E_2, \dots, E_n$ be any $n$ mutually exclusive events of a random experiment such that $P(E_i) \ne 0$ for all $i=1, 2, \dots, n$, and let $A$ be any arbitrary event which is a subset of $\bigcup_{i=1}^n E_i$, then calculate $P(E_i|A)$ for all $i=1, 2, \dots, n$ in the case when $P(A)>0$.
  \item There are two bags A and B. A contains $n$ white and 2 black balls, and B contains 2 white and $n$ black balls. One of the two bags is selected at random and 2 balls are drawn from it without replacement. If both the drawn balls are white and the probability that the bag A was used to draw the balls is $\frac{6}{7}$, then find the value of $n$.
\end{parts}
\medskip\rule{\linewidth}{0.2pt}

\medskip\noindent\textbf{Question 6.}
\begin{parts}
  \item Define mathematical expectation of a random variable. Further, state and prove the multiplication theorem of expectation for any two independent random variables $X$ and $Y$.
\end{parts}
\medskip\rule{\linewidth}{0.2pt}

\medskip\noindent\textbf{Question 7.}
\begin{parts}
  \item Define probability density function (p.d.f.) of a continuous random variable. Next, if $X$ denotes the diameter of an electric cable and has the following probability density function:
    $$f(x) = 6x(1-x), \quad 0 < x < 1$$
    and $0$ otherwise. Then:
    \begin{parts}
      \item check whether $f(x)$ is a p.d.f., and
      \item determine a number $b$ such that $P(X < b) = P(X > b)$.
    \end{parts}
  \item If the joint p.d.f. of $X$ and $Y$ is:
    $$f(x, y) = \frac{8}{9}xy, \quad 1 \le x \le y \le 2$$
    and $0$ otherwise, then find the conditional p.d.f. of $Y$ given $X = x$ and that of $X$ given $Y = y$.
  \item Define joint p.m.f. and joint p.d.f. Next, the joint probability distribution of two random variables $X$ and $Y$ is given by:
    $$P(X=0, Y=1) = \frac{1}{4}, \quad P(X=1, Y=-1) = \frac{1}{2}, \quad P(X=1, Y=1) = \frac{1}{4}$$
    Then find:
    \begin{parts}
      \item the marginal probability distributions of $X$ and $Y$, and
      \item the conditional probability distribution of $X$ given $Y=y$.
    \end{parts}
\end{parts}
\medskip\rule{\linewidth}{0.2pt}

\medskip\noindent\textbf{Question 8.}
\begin{parts}
  \item Derive the expression for the variance of the linear combination of any $n$ random variables.
  \item Calculate the expectation of the number of failures preceding the first success in an infinite series of independent trials with constant probability of success, say $p$, in each trial.
\end{parts}

\vfill
\rule{\linewidth}{0.4pt}
\begin{center}\small
  \href{https://bhu-exam.vercel.app/}{Click here to visit our website}\\[4pt]\href{https://me-aryan.vercel.app/}{Click here to visit our contributor}\\[4pt]\href{https://quashan-me.vercel.app/}{Click here to visit our contributor}
\end{center}

\end{document}